Finite volume methods blend well with laser heat deposition as the line heat source is easily integrated over each cell the laser passes through. The total energy deposited in any cell is proportional to the amount
of intensity attenuated as the laser passes through that cell.

The coordinate system is local to each droplet such that the Cartesian coordinate axes align with the principal axes, and the ellipsoidal center is the coordinate origin. To ease calculations, the position and direction
vectors, $\mathbf{r}$ and $\mathbf{\hat{\upomega}}$ respectively, are linearly mapped into the droplet's basis vectors. Upon entrance to a cell $(i,j,k)$ at position $\mathbf{r}$, the exit surface on that cell is one with the
smallest positive travel distance. Computing the distances is therefore the main routine of source term calculations.

Surfaces of constant $r$ are ellipsoids, and the distance to any such surface is found using Equation \ref{eqn:distance_ellipsoid}. Again, the ellipsoidal center $\mathbf{r}_0$ is the origin, and the right-hand side
becomes $r^2$.

Surfaces of constant $\mu$ are elliptic cones of the form
\begin{equation}
    (1-\mu^2) \frac{x'^2}{a^2} - \mu^2 \left( \frac{y'^2}{b^2} + \frac{z'^2}{c^2} \right) = 0,
\end{equation}
or in matrix notation
\begin{subequations}
    \label{eqn:heat_ellipsoidal_mu}
    \begin{align}
        \mathbf{r} &\cdot \mathbf{M}\mathbf{r} = 0 \\
        \mathbf{M} &\equiv \diag
            \begin{bmatrix}
                \displaystyle \frac{1-\mu^2}{a^2} & \displaystyle -\frac{\mu^2}{b^2} & \displaystyle -\frac{\mu^2}{c^2}
            \end{bmatrix}
    \end{align}
\end{subequations}

Substituting $\mathbf{r} := \mathbf{r} + \mathbf{\hat{\upomega}s}$ into Equation \ref{eqn:heat_ellipsoidal_mu} yields a quadratic equation in $s$. It's worth noting that the cones described in Equation
\ref{eqn:heat_ellipsoidal_mu} include surfaces $\mu$ and $-\mu$, so it is important to ensure that $\mathbf{r} + \mathbf{\hat{\upomega}s}$ lands on the correct surface.

Surfaces of constant $\varphi$ are ellipses confined to the plane
\begin{equation}
    \frac{y'}{b} \sin\varphi = \frac{z'}{c} \cos\varphi,
    \label{eqn:heat_ellipsoidal_phi}
\end{equation}
which will yield a linear equation in $s$ after substituting $y' := y' + \omega_y' s$ and $z' := z' + \omega_z' s$. Again, after solving for $s$ it is important to ensure the correct surface is
landed on, as Equation \ref{eqn:heat_ellipsoidal_phi} is identical for $\varphi$ and $\varphi \pm \pi$.

% \begin{equation}
%     \left( \mathbf{r} + \mathbf{\hat{\omega}}s \right) \cdot \mathbf{M} \left(\mathbf{r} + \mathbf{\hat{\omega}}s \right) = r^2,
%     \label{eqn:heat_distance_r_matrix}
% \end{equation}
% which is a quadratic equation of the form
% \begin{subequations}
%     \begin{equation}
%         As^2 + B^2 + C = 0
%     \end{equation}
%     where the coefficients $A$, $B$ and $C$ are defined as
%     \begin{align}
%         A &= \mathbf{\hat{\omega}} \cdot \mathbf{M}\mathbf{\hat{\omega}} \\
%         B &= \mathbf{\hat{\omega}} \cdot \mathbf{M}\mathbf{r} \\
%         C &= \mathbf{r} \cdot \mathbf{M} \mathbf{r} - r^2.
%     \end{align}
% \end{subequations}

The routine terminates once the ray exits surface $r = 1$.